\documentclass[10pt,a4paper]{amsart}
\usepackage[margin=19mm]{geometry}
\usepackage{amsmath,amssymb,mathtools}
\usepackage[scr=rsfs,cal=euler]{mathalfa}
\usepackage{xfrac}
\usepackage[unicode,hidelinks,bookmarks=false]{hyperref}
\newcommand{\ie}{i.e.\ }
\newcommand{\cf}{cf.\ }
\newcommand{\resp}{respectively\ }
\newcommand{\Subsection}{\subsection}
\providecommand{\nobreakdash}{\nobreak\hbox{-}}
\setlength{\emergencystretch}{2em}
\multlinegap0pt
\usepackage{mathtools}
\usepackage[shortlabels]{enumitem}
\usepackage{amssymb}
\usepackage{csquotes}
\newtheorem{lemma}{Lemma}
\newcommand{\N}{\mathbb{N}}
\newcommand{\R}{\mathbb{R}}
\newcommand{\lm}{\mathscr{L}}
\title{Ultralimits of Sobolev maps and stability of Dehn functions}
\author{Toni Ikonen, Stefan Wenger}
\date{Source: arXiv:2603.05246v1 (CC BY 4.0)}
\begin{document}
\maketitle

\noindent\emph{Source and licence.} Ultralimits of Sobolev maps and stability of Dehn functions. Version arXiv:2603.05246v1 (CC BY 4.0), distributed by arXiv under the Creative Commons Attribution 4.0 International licence, http://creativecommons.org/licenses/by/4.0/, as recorded in the arXiv licence metadata for this exact version and shown on the abstract page https://arxiv.org/abs/2603.05246v1. Full licence notice: \texttt{licenses/arxiv-2603.05246-CC-BY-4.0.txt}.\par\smallskip

\noindent\textit{Editorial scope.} Counted unit: the article's lemma lem:bound-ultralimit-set (source0.tex lines 453--455). The article's complete original proof of it is delivered with it (source0.tex lines 457--481, one \texttt{proof} environment of 25 lines). No mathematics is written, repaired or re-derived: the statement and the proof are the article's own lines, in the article's own notation, and no retained line is substituted or re-flowed.  No further source-local context is needed: every cross-reference of every retained line resolves inside the two delivered spans.  Nothing was omitted from any retained span, so the package carries no omission record.  No retained line contains a citation, so the wrapper carries no bibliography and no bibliography entry.  No retained line contains a figure environment, an image-inclusion command or drawing code, so no image is delivered and the package declares no asset dependency.  Editorial formatting only, in addition: an AMS \texttt{amsart} preamble that restates the article's own declarations and macros used by the retained lines, namely the environments \texttt{lemma} on separate counters, together with the standard packages those lines need.  The retained environments print in this document as Lemma 1 in order of appearance, whereas the article prints the same items with the numbering of its own full document.  The complete native source of the article as distributed in the arXiv e-print for this exact version is bundled unchanged as \texttt{sources/195787/source0.tex}.
\begin{lemma}\label{lem:bound-ultralimit-set}
Let $(E_m)$ be a sequence of non-empty measurable subsets of a bounded open set $\Omega \subset \R^n$. Then the ultralimit $E_\omega$ is a non-empty compact subset of $\overline{\Omega}$ and satisfies $$\lm^n(\Omega\cap E_\omega) \geq \lim\nolimits_\omega \lm^n(E_m).$$
\end{lemma}

\begin{proof}
Clearly $E_\omega$ is a non-empty compact subset of $\overline{\Omega}$ so it suffices to prove the measure lower bound for $\Omega \cap E_\omega$.
%
%The boundedness of $\Omega$ implies that $E_\omega$ is not empty. The closedness is standard so compactness follows by boundedness.\todo{S: I wonder whether it wouldn't be enough to simply say "Clearly, $E_\omega$ is a non-empty compact subset of $\overline{\Omega}$ and it therefore suffices to prove the lower bound on the measure of $\Omega\subset E_\omega$ when $\lim\nolimits_{\omega} \lm^n( E_m ) > 0$.}
%
%It remains to consider the measure bound when $\lim\nolimits_{\omega} \lm^n( E_m ) > 0$. 
Observe that the sequence of characteristic functions $g_m = \chi_{ E_m }$ forms a bounded sequence in $L^{\infty}( \mathbb{R}^n )$ and as $L^{\infty}( \mathbb{R}^n )$ is the dual of $L^{1}( \mathbb{R}^n )$, there exists $g_\omega \in L^{\infty}( \mathbb{R}^n )$ such that
\begin{align*}
    \int_{ \mathbb{R}^n } h g_\omega \,dz = \lim\nolimits_{\omega} \int_{ \mathbb{R}^n } h g_m \,dz
    \quad\text{for every $h \in L^{1}( \mathbb{R}^n )$.}
\end{align*}
In particular, if we apply this equality to nonnegative $h$, it immediately follows that $0 \leq g_\omega \leq 1$ almost everywhere. We conclude that
\begin{align*}
    \lm^n( \Omega \cap \left\{ g_\omega > 0 \right\} )
    \geq
    \lim\nolimits_{\omega} \int_{ \Omega } g_m \,dz
    =
    \lim\nolimits_{\omega} \lm^n( E_m ).
\end{align*}
If $x \in \Omega \cap \left\{ g_\omega > 0 \right\}$ is a Lebesgue density point, we claim, in fact, that $x \in E_\omega$. This will finish the proof. Let
\begin{align*}
    \mathcal{G}_k = \{ m \in \N \colon \lm^{n}( E_m \cap B( x, k^{-1} ) ) > 0 \} \cap [k,\infty).
\end{align*}
It holds that $\mathcal{G}_{k} \supset \mathcal{G}_{k+1}$ and $\omega( \mathcal{G}_k ) = 1$ for $k \in \N$ by construction. To finish, let $x_m \in E_m$ for $m \in \N \setminus \mathcal{G}_1$ and $x_m \in E_m \cap B( x, k^{-1} )$ if $m \in \mathcal{G}_{k} \setminus \mathcal{G}_{k+1}$. Since $x = \lim\nolimits_{\omega} x_m$, the proof is complete.
\end{proof}

\end{document}
